\begin{proposicion}[Concentration of the observables]\label{prop:concentracion-meso}
For a fixed family of patterns and fiber types with bounded weights, fix $b_1>0$ and $\chi>0$.
There is a rate $c_{\rm pat}>0$ such that, for each floor $K$, the following event has probability
at most $C e^{-c_{\rm pat}K}$ under $\nu$. The history enters regularly at $M$ and some observable
at a reached size $k\ge K$, up to and including the first-exit board, has deviation greater than
\[
 C_H\{A(\delta+1/n)+\chi+1/K\}+b_1 .
\]
We measure deviations from the value at $M$, comparing each surviving tracked line with its own fiber.
The rate depends on the constants of the outer class, the family of observables, $b_1$, and $\chi$.
\end{proposicion}
\begin{proof}
We subtract from each increment its conditional expectation, which is predictable, and freeze the
process upon reaching the first-exit board. For a pattern with no tracked line, the centered jump in
the step from size $u$ is at most $2C_H/u$. Azuma's maximal inequality \eqref{macro:max-mart}, with
$\sum_{u>k}u^{-2}\le1/k$, bounds by $2\exp(-c_H b_1^2k)$ the probability that the absolute value of
the martingale reaches $b_1$ between $M$ and $k$.
Lemma~\ref{lema:presupuesto-meso} and the sum of $u^{-2}$ bound the cumulative drift simultaneously for all
prefixes.

Fix a history of regular entry, leaving the unrevealed row order random, and fix one of its tracked
lines. At a regular parent of size $u$, restrict the original transition to children where the line
survives and divide by their total probability $y_u$. This conditions just the next step. Stop this
auxiliary law at the first exit or at the floor. On each surviving prefix from $M$ to $k$ its
transition probabilities are the original ones divided by $y_u$ at each successive parent. Hence the
original prefix weight given this entry history equals the auxiliary prefix weight multiplied by
\[
 \prod_{u=k+1}^{M}y_u
 =\prod_{u=k+1}^{M}\big(1-1/u+O(u^{-2})\big)
 =(k/M)\exp(O(1/k)).
\]
For a row the product is exactly $k/M$. For a column, dividing each factor by $1-1/u$ and summing
the logarithms gives an error $O(\sum_{u>k}u^{-2})=O(1/k)$, uniformly in the prefix. The same weight
identity applies prefix by prefix at a stopping time, including a step that first exits. Under the
auxiliary kernel, we subtract the conditional drift of
Lemma~\ref{lema:incrementos-meso} and freeze the process upon reaching the first-exit board. Azuma's
maximal inequality holds in the same way starting from any history with regular entry.

In the band $m\le k<2m$, with $m=2^jK$, stop at the first reached board in the band at which the
martingale for this tracked line has absolute value $\ge b_1$, if there is one. These prefixes are
disjoint, and their auxiliary probability is bounded by the maximal inequality down to size $m$.
Transferring them to $\nu$ introduces a factor of at most $(2m/M)e^{C/m}$. A union bound over the $M$
tracked lines of each family and the fixed weight types, followed by a sum over the bands, gives at most
$\sum_m C m e^{-c_H b_1^2 m}\le C'e^{-c' b_1^2K}$. As long as the tracked line remains available,
the drift subtracted along these trajectories is exactly that of the auxiliary kernel. The diagonal
budget along the original trajectory controls the accumulated absolute value of this drift. After
integrating over the entry history, we count failure of the budget just once under the original law
$\nu$, using the \emph{first} estimate of Lemma~\ref{lema:presupuesto-meso}. Take $c_{\rm pat}$ smaller
than both its rate $c_\chi$ and all the concentration rates above, after summing the bands. The
geometric rate $\lambda$ in the second estimate of that lemma is chosen separately.
\end{proof}

\subsection{Reconstruction of the regularity conditions}\label{meso:regularidad}
Write $W_0=\nu_1/(u_Du_A)=b_0r_Rr_C$, the value at $x=0$ of the weight $W_x$ from
Proposition~\ref{macro:trayectoria}, and set $q_{ij}=Y_{ij}W_{0,ij}/k$ and $V_{ij}=\nu_{1,ij}/k$, evaluated
at the original coordinates of the square. The density bounds $c_\nu,C_\nu$, the minimum vacancy
$\kappa_\nu$, and the Lipschitz constant $L_\nu$ are the fixed profile constants. Consider a fixed partition
of $[0,1]$ into intervals of length at most $a_{\rm mes}$. The radius $\varrho_{\rm mac}$ and mesh $a_{\rm
mac}$ of Section~\ref{macro} are chosen later, in item~5 of Section~\ref{meso:constantes}, as functions of
this partition, the tolerances below, and $\delta$ (Lemma~\ref{meso:entrada}). For rows and columns we
control the proportion of the remaining indices in each interval, with the interval length as target. Fix a
tolerance $\varrho_{\rm mes}>0$ and individual tolerances whose sum is at most $\varrho_{\rm mes}$ per
family, for example $\varrho_{\rm mes}|I|$ for $I$ in this partition. Each row or column fiber has weights
$W_0$, target $1$, and tolerance at most $\varrho_{\rm mes}$.

To describe the other patterns, write $T=\{1,2\}\times\{1,2\}$. For each $S\subseteq T$, take the pattern on
two rows and two columns that requires an available square at the positions in $S$. Its weight is
$\prod_{e\in S}W_{0,e}\prod_{e\in T\setminus S}\nu_{1,e}$. All sixteen patterns have the same target,
\[
 \int_{[0,1]^4}\nu_1(r,c)\nu_1(r,c')\nu_1(r',c)\nu_1(r',c')\,dr\,dr'\,dc\,dc' .
\]
The normalization for injective assignments is $(k)_2^2$. Using $k^4$ and including assignments that repeat
a row or a column changes each pattern by $O(1/k)$, since the $O(k^3)$ extra assignments have bounded
weight. With $E=q-V$, as in Lemma~\ref{macro:brecha}, the identity
\[
 \operatorname{tr}\big[(EE^{T})^2\big]
 =\sum_{i,i',j,j'} E_{ij}E_{i'j}E_{i'j'}E_{ij'}
\]
produces the sixteen terms extended in this way with signs $(-1)^{4-|S|}$. The sixteen targets cancel since
$(1-1)^4=0$. Choose the positive tolerances individually so that their sum, plus $O(1/K_0)$, is smaller than
the fourth power of the required spectral bound. Then $\|q-V\|_{\rm op}\le\operatorname{tr}[(EE^T)^2]^{1/4}$
is smaller than that bound. The expansion includes $S=\varnothing$ and permits repeated diagonals in every
term.
